\chapter*{How to read this book}
\addcontentsline{toc}{chapter}{How to read this book}
\markboth{How to read this book}{How to read this book}

The book is five parts and an epilogue, and they build. Part~I is load-bearing
for everything after it. The rest can be raided.

\begin{description}[leftmargin=0pt, style=nextline, itemsep=8pt]

\item[Part I, Protocol Fundamentals (Chapters 1 to 4)]
The physics, then the wire. Chapter~\ref{ch:primer} sets up the three budgets
every agentic application spends: latency, context, and cost. Every later
chapter cashes out against those three. Chapters~\ref{ch:architecture}
to~\ref{ch:authorization} are the stateless core, both transports down to the
bytes, and OAuth as MCP actually uses it. Do not skip Chapter~\ref{ch:primer}.
It is short and the rest of the book assumes it.

\item[Part II, The Primitives (Chapters 5 to 10)]
Tools, resources, prompts, multi round-trip requests, the extensions framework
with Tasks, and MCP Apps. Every chapter has the same shape: how it works on the
wire, what it costs, and when to use it. These are reasonably independent, so
read the ones you need.

\item[Part III, Building MCP Applications (Chapters 11 to 14)]
The hands-on arc. Servers, hosts, the agent loop, and multi-agent topologies.
If you learn by building, start here and jump backwards when something does not
make sense.

\item[Part IV, Performance Engineering (Chapters 15 to 18)]
Test it, measure it, optimise it, run it. This is the part the book exists for.
Chapter~\ref{ch:optimizing} is the one to photocopy and pin up.

\item[Part V, Security and Trust (Chapters 19 and 20)]
Threat modelling and access control. Read Chapter~\ref{ch:threats} before you
expose anything to the public internet, not after.

\end{description}

\section*{Three routes through}

\textbf{You are building a server.} Chapters~\ref{ch:primer},
\ref{ch:architecture}, \ref{ch:transport}, \ref{ch:tools}, \ref{ch:resources},
\ref{ch:building-servers}, \ref{ch:testing}, then \ref{ch:threats}. Then
Chapter~\ref{ch:optimizing} once it works.

\textbf{You are building a host or a client.} Chapters~\ref{ch:primer},
\ref{ch:architecture}, \ref{ch:transport}, \ref{ch:mrtr},
\ref{ch:building-hosts}, \ref{ch:agent-loop}, then all of Part~IV. The host is
the trust boundary, so Chapter~\ref{ch:access-control} matters more to you than
to anyone else.

\textbf{Something is slow and you need it fixed by Friday.}
Chapter~\ref{ch:primer} for the budgets, Chapter~\ref{ch:measuring} to find out
where the time goes, Chapter~\ref{ch:optimizing} to get it back.
Appendix~\ref{ch:perf-checklist} is the one-page version if it needs fixing by
Wednesday.

\section*{The furniture}

Four kinds of box appear throughout, and they mean specific things.

\begin{legacybox}{What a legacy box looks like}
A feature that is Deprecated under the lifecycle policy, or a pattern from an
earlier revision. It is here so you can recognise and maintain it. It is never
a recommendation.
\end{legacybox}

\begin{measurebox}{What a measurement box looks like}
A number from an actual run of the Meridian harness. If part of it is modelled
rather than measured, the box says which part.
\end{measurebox}

\begin{dangerbox}{What a threat box looks like}
An attack that works. These appear from Chapter~\ref{ch:architecture} onwards,
not just in Part~V, because security is a design input and not a final coat of
paint.
\end{dangerbox}

\keyidea{And a line like this is the one-sentence version of the argument you
just read. If you remember nothing else from a section, remember these.}

Code listings come in three flavours. Wire dumps on a blue ground are literal
JSON-RPC, exactly as it appears on the connection. HTTP exchanges show real
headers, including the ones \texttt{2026-07-28} made mandatory. Everything else
is Python or TypeScript, extracted from the Meridian repository, which means it
compiles and its tests pass.

\section*{The companion repository}

\begin{center}
\texttt{https://github.com/krimler/Building-Awesome-MCP-Apps}
\end{center}

\noindent Clone it before Chapter~\ref{ch:primer} if you want to follow along.
You need Python 3.11 or later and nothing else at all.

\begin{sh}
git clone https://github.com/krimler/Building-Awesome-MCP-Apps
cd Building-Awesome-MCP-Apps
make test     # 191 tests, about two seconds
make bench    # regenerates every number printed in this book
\end{sh}

\noindent Appendix~\ref{ch:companion} is the guided tour.

\cleardoublepage
